% =============================================================================
%  Chapter 5 -- Conclusion
% =============================================================================
\chapter{Conclusion}
\label{sec:conclusion}

% STYLE: four moves, and nothing else.
%
%   1. What you did and what you found -- a compressed restatement, in the past
%      tense. No new results may appear here for the first time.
%   2. THE ANSWER to the research question posed in \Cref{sec:research_question}.
%      See below -- this is not optional, and not the same as move 1.
%   3. Limitations -- honestly scoped. What the results do NOT license anyone to
%      conclude.
%   4. Future work -- specific enough that someone could pick one up and start.
%      "More data" is not future work; "extending the evaluation to X, where the
%      assumption of Y no longer holds" is.
%
% Length: two to four pages. A conclusion that restates the introduction at
% length is a conclusion nobody reads.

Placeholder summary of what was done and found.

% ANSWERING THE RESEARCH QUESTION. Summarising the work is not the same as
% answering the question, and a summary is what most drafts stop at. Restate
% the question, then answer it in a sentence or two, in plain language, and say
% what the answer rests on.
%
% Two rules:
%   * Answer the question you actually asked in \Cref{sec:research_question}.
%     If the thesis ended up answering a different one, go back and change the
%     question -- do not quietly substitute it here.
%   * A qualified or negative answer is a real answer. "No, and here is the
%     condition under which it fails" is a finding. Hedging until no answer is
%     discernible is not.
%
% If sub-questions were posed, answer each, then the main question.

Returning to the research question posed in \Cref{sec:research_question}:
placeholder answer, stated plainly and attributed to the evidence that
supports it.

Placeholder limitations and future work.
